\chapter{数式処理システム入門}\label{cas}

数式処理システムとは，シンボリックな計算をサポートするプログラミング言語のことをいう．
数式処理システムは，未束縛の変数をシンボルとして扱う．
数式処理システムを使うと，たとえば，$x + x \rightarrow 2x$や$(x + y)^2 \rightarrow x^2 + 2 x y + y^2$のような計算ができる．
Egisonもこのようなシンボリックな計算をサポートしている．
本章は数式処理システムを使うとどのような計算ができるのかみていく．

\section{シンボルの宣言}\label{cas-symbol}

数式処理システムは，未束縛の変数をシンボルとして扱う．
Egisonでは，シンボルを使う前に\lstinline{declare symbol}文で宣言する必要がある．
シンボルには型を指定でき，省略した場合は\lstinline{MathExpr}型が使われる．
\begin{lstlisting}[language=egison,numbers=none]
declare symbol x, y, z : MathExpr

x -- x
\end{lstlisting}

数式処理システムとしての機能をもつEgisonには，シンボル同士の足し算や掛け算が組み込みで定義されている．
そのため，たとえば，$x + x \rightarrow 2x$や$(x + y)^2 \rightarrow x^2 + 2 x y + y^2$のような計算ができる．
\begin{lstlisting}[language=egison,numbers=none]
x + x -- 2x
(x + y)^2 -- x^2 + 2 * x * y + y^2
\end{lstlisting}

Egisonは数式を自動で積和標準形に展開する．
積和標準形とは，掛け算が内側に，足し算が外側になるような形の数式のことである．
そのため，$(x + y)^2$は$x^2 + 2 x y + y^2$のように展開される．

出力の数式の項の因子として現れるシンボルの順番は，ユーザーが入力した数式に現れる順番と同じになるように実装されている．
たとえば，上記の$(x + y)^2$の$x$と$y$を入れ替えて$(y + x)^2$を計算すると$2$項目は$2 y x$となる．
\begin{lstlisting}[language=egison,numbers=none]
(y + x)^2 -- y^2 + 2 * y * x + x^2
\end{lstlisting}

\section{関数適用のシンボル化}\label{cas-symbol2}

\lstinline{sqrt x}や\lstinline{exp x}のようにシンボルを引数にとる可能性がある関数について，関数適用を止めて，その式をシンボリックな値としてそのまま返すようにライブラリで定義されている．
\begin{lstlisting}[language=egison,numbers=none]
declare symbol x : MathExpr

sqrt 4 -- 2
sqrt x -- sqrt x
exp 1 -- e
exp x -- exp x
\end{lstlisting}

\lstinline{sqrt 2}のように評価を進められない関数適用についても，関数適用を止めて，その式をそのまま返すようにライブラリで定義されている．
$\sqrt{8} \rightarrow 2 \sqrt{2}$のような簡約もライブラリで定義されている．
\begin{lstlisting}[language=egison,numbers=none]
sqrt 2 -- sqrt 2
sqrt 8 -- 2 * sqrt 2
\end{lstlisting}
これらの定義については，\ref{mexpr-squote}節で紹介するシングルクオートを使って定義される．


\section{数式の簡約}\label{cas-simplify}

$\sqrt{x} \cdot \sqrt{x} = x$や$sin^2\theta + cos^2\theta = 1$のような数式を紙の上で扱うときにおこなう簡約が，Egisonを含む多くの数式処理システムには実装されている．
\begin{lstlisting}[language=egison,numbers=none]
declare symbol x, `$\theta$` : MathExpr

sqrt x * sqrt x -- x
(sin `$\theta$`)^2 + (cos `$\theta$`)^2 -- 1
\end{lstlisting}
Egisonでは，このような簡約のための機能は，Sweet Egison(Egisonのパターンマッチ機能を提供するHaskellライブラリ)を使ってHaskellにより言語に組み込みで実装されている．

上記以外にも，複素数などについての簡約ルールも実装されている．
\begin{lstlisting}[language=egison,numbers=none]
i^2 -- -1
w^3 -- 1  
w + w^2 -- -1  
\end{lstlisting}
簡約ルールは\url{https://github.com/egison/egison/blob/master/hs-src/Language/Egison/Math/Normalize.hs}で定義されている．

数式の簡約の自動化は非常に難しく，Egisonでは簡潔なアルゴリズムしか用意されていない．
たとえば，多項式の約分にはグレブナー基底についての理論に基づいたアルゴリズムが有用であるが，Egisonには実装されていない．
そのため，Egisonは複雑な多項式の分数の適切に簡約できないことがある．
なお，単変数の多項式については，$(x^2-1)/(x-1) \rightarrow x+1$のようなユークリッドの互除法による約分が実装されている．

Wolfram言語(Mathematica)は，このような高度な数式の簡約ができる．
``\texttt{-M mathematica}''オプションを追加してEgison処理系を実行するとWolfram言語で読み込めるかたちで数式を出力する．
この出力をWolfram言語の処理系に入力すれば，Wolfram言語の高度な簡約機能を利用できる．
\begin{lstlisting}[language=egison,numbers=none]
$ egison -M mathematica
> x + sqrt 2
#mathematica|x + Sqrt[2]|#
\end{lstlisting}

\section{簡約規則の宣言}\label{cas-rule}

前節で紹介した組み込みの簡約規則に加えて，ユーザーが\lstinline{declare rule}文で簡約規則を宣言することもできる．
たとえば，$\sqrt{2}$を表すシンボル\lstinline{sqrt2}を導入し，$\sqrt{2}^2 = 2$という規則を登録すると，この規則は簡約のたびに自動で適用される．
\begin{lstlisting}[language=egison,numbers=none]
declare symbol sqrt2
declare rule auto term sqrt2^2 = 2

(1 + sqrt2)^2 -- 2 * sqrt2 + 3
\end{lstlisting}
これは数学的には$\mathbb{Z}[\sqrt{2}] = \mathbb{Z}[X]/(X^2-2)$という代数拡大（多項式環の商）を導入する操作である．
Egisonの簡約規則は大域的であるが，シンボルは大域的に一意であるため，規則はそのシンボルが現れる数式にしか作用しない．
つまり，シンボルの一意性がスコープの役割を果たしている．
ガウス整数$\mathbb{Z}[i]$（$i^2 = -1$）や二重数（$\varepsilon^2 = 0$），1の冪根なども，同じパターンで導入されている代数拡大である．

\section{型による正規形の指定}\label{cas-tower}

同じ数式は，複数の異なるデータ構造で表現できる．
たとえば，
\[ (2+3i) + (1-i)x + 4x^2 \]
は，原子$\{i, x\}$の上の平坦な多項式としても，係数がガウス整数$\mathbb{Z}[i]$であるような$x$の多項式としても保持できる．
どちらの正規形が適切かは，つぎに何をしたいか（$x$について微分したいのか，$\mathbb{Z}[i]$の性質を使いたいのか）に依存するため，値だけからは決められない．
Egisonでは，\textbf{型注釈でこの正規形を指定する}．
\begin{lstlisting}[language=egison,numbers=none]
declare symbol i, x

def v := (2 + 3 * i) + (1 - i) * x + 4 * x^2

def flat : Poly Integer [i, x] := v
inspect flat
-- "4 * x^2 - i x + x + 3 * i + 2 : Poly Integer [i, x]"

def byX : Poly (Poly Integer [i]) [x] := v
inspect byX
-- "4 * x^2 + (- i + 1) * x + 3 * i + 2 : Poly (Poly Integer [i]) [x]"
\end{lstlisting}
型\lstinline{Poly Integer [i, x]}は「原子集合$\{i, x\}$の上の整数係数多項式」を表す型である．
\lstinline{Poly (Poly Integer [i]) [x]}のように係数の位置に多項式型を入れ子にすると，「$\mathbb{Z}[i]$係数の$x$の多項式」という正規形を表す．
型注釈は値を分類するだけでなく，値の内部構造をその型のかたちに組み替える（reshape）．
組み替えは値を変えない．
実際，2つの正規形の差を計算すると$0$になる．
\begin{lstlisting}[language=egison,numbers=none]
flat - byX -- 0
\end{lstlisting}
形の異なる表現同士の演算が正しく簡約されるのは，\textbf{演算の結果は常に既定の平坦形で返る}という規約があるためである．
入れ子の正規形は型注釈を書いた場所にだけ現れ，演算をまたいで勝手に伝播することはない．

上の例で使った\lstinline{inspect}は，値と一緒にその値から読み取れる型（観察型）を表示する．
よく使う型には\lstinline{declare cas-type}文で別名をつけられる．
別名は透明であり，型の表示では展開されたもとの型が使われる．
\begin{lstlisting}[language=egison,numbers=none]
declare cas-type GaussianPolyX := Poly (Poly Integer [i]) [x]

def byX2 : GaussianPolyX := v
inspect byX2
-- "4 * x^2 + (- i + 1) * x + 3 * i + 2 : Poly (Poly Integer [i]) [x]"
\end{lstlisting}

\section{部分型関係と暗黙の昇格}\label{cas-subtype}

CASの型は，値を保存する埋め込みによって順序づけられている．
たとえば，整数は有理数に埋め込むことができ（\lstinline{Integer} $\subset$ \lstinline{Frac Integer}），有理数は有理数係数の多項式に定数として埋め込むことができる．
異なる型の値を混ぜて演算すると，両方を含む最小の型（join）へ自動で昇格する．
\begin{lstlisting}[language=egison,numbers=none]
def p : Poly Integer [x] := x + 1
def q : Frac Integer := 1 / 2

def r := p + q
inspect r -- "x + 3 / 2 : Poly (Frac Integer) [x]"
\end{lstlisting}
原子集合が異なる場合も，原子集合を合併した型へ昇格する．
\begin{lstlisting}[language=egison,numbers=none]
declare cas-type GaussianInt := Poly Integer [i]
declare cas-type Zsqrt2 := Poly Integer [sqrt2]

def a : GaussianInt := 1 + i
def b : Zsqrt2 := 1 + sqrt2

def c := a + b
inspect c -- "sqrt2 + i + 2 : Poly Integer [i, sqrt2]"
\end{lstlisting}
$\mathbb{Z}[i]$の値と$\mathbb{Z}[\sqrt{2}]$の値の和が，注釈なしで$\mathbb{Z}[i, \sqrt{2}]$に着地している．

組み込みの順序は，前節の平坦形（\lstinline{Poly Integer [i, x]}）と入れ子形（\lstinline{Poly (Poly Integer [i]) [x]}）をあえて関係づけていない．
両者は同じ値の集合の異なる正規形であり，どちらへ昇格すべきかはユーザーの選択だからである．
この関係は\lstinline{declare cas-subtype}文でユーザーが宣言できる．
\begin{lstlisting}[language=egison,numbers=none]
declare cas-subtype Poly Integer [i, x] <: Poly (Poly Integer [i]) [x]
\end{lstlisting}
辺を宣言するたびに，順序全体で任意の2つの型のjoinが一意に定まること（join半束性）が検査される．
一意性を壊す宣言は，先に宣言すべき辺の提案つきでエラーになるため，順序は常に健全に保たれる．

\section{商の型}\label{cas-quotient}

$\mathbb{Z}/7\mathbb{Z}$（7で割った余りの世界）のような商は，これまでの型とは性質が異なる．
$12$と$5$は$\mathbb{Z}$では異なる値だが，$\mathbb{Z}/7\mathbb{Z}$では等しい．
つまり\textbf{等しさそのものが型に依存する}ため，値を保存する埋め込みの順序（前節）には入れられない．
Egisonでは，商の型は\lstinline{declare cas-quotient}文で順序の外に宣言する．
\begin{lstlisting}[language=egison,numbers=none]
declare cas-quotient Mod7 := Integer by (\n -> modulo n 7)

def a : Mod7 := projMod7 5
def b : Mod7 := projMod7 4

a + b -- 2
projMod7 12 == projMod7 5 -- True
\end{lstlisting}
宣言すると，整数から\lstinline{Mod7}への射影\lstinline{projMod7}，代表元を取り出す\lstinline{reprMod7}，環としての演算，等値判定が自動で導出される．
演算は毎回代表元を簡約するので，値は常に$0$から$6$の範囲に保たれる．
また宣言時には，簡約関数が冪等であることと，加算・乗算と両立する（合同である）ことが，サンプル点の上で検査される．

\lstinline{Mod7}の値と生の整数の混在（\lstinline{a + 10}など）は型エラーになる．
これは$\mathbb{Z}$の等しさと$\mathbb{Z}/7\mathbb{Z}$の等しさを暗黙に混同しないためであり，世界を横断するときは\lstinline{projMod7}・\lstinline{reprMod7}を明示的に使う．
商の上に直接関数を定義することもできる．
たとえば$\mathbb{F}_7$における逆元は，フェルマーの小定理（$z^{-1} = z^5$）を使って1行で書ける．
\begin{lstlisting}[language=egison,numbers=none]
def inv7 (z : Mod7) : Mod7 := z * z * z * z * z

reprMod7 (inv7 (projMod7 3)) -- 5
\end{lstlisting}

\section{二次方程式を解くプログラム - 数式を処理するアルゴリズムの記述}\label{cas-quadratic}

本節では，二次方程式を解くアルゴリズムをプログラムとして書いていく．
以下のように二次方程式を引数にとり，解を返す関数\lstinline{qF}を定義していく．
\begin{lstlisting}[language=egison,numbers=none]
qF (x ^ 2 + x + 1) x
-- ((-1 + i * sqrt 3) / 2, (-1 - i * sqrt 3) / 2)

qF (x ^ 2 + b * x + c) x
-- ((- b + sqrt (b^2 - 4 * c)) / 2, (- b - sqrt (b^2 - 4 * c)) / 2)

qF (a * x ^ 2 + b * x + c) x
-- ((- b + sqrt (b^2 - 4 * a * c)) / (2 * a), (- b - sqrt (b^2 - 4 * a * c)) / (2 * a))

qF (a * x ^ 2 + 2 * b * x + c) x
-- ((- b + sqrt (b^2 - a * c)) / a, (- b - sqrt (b^2 - a * c)) / a)
\end{lstlisting}
\lstinline{qF}は解の公式$\frac{-b \pm \sqrt{b^2 - 4 a c}}{2 a}$をそのまま書くだけでも定義することができる．
以下のプログラムは\lstinline{qF'}という補助関数を使って\lstinline{qF}を定義している．
\lstinline{qF'}は第一引数に二次の係数，第二引数に一次の係数，第三引数に定数項をとる．
たとえば，\lstinline{qF' 1 0 2}は2次方程式$x^2 + 2 = 0$の解を返す．
\lstinline{coefficients}は多項式とシンボルを引数にとり，そのシンボルについて係数のリストを返す関数である．
\lstinline{qF'}は単純に$\frac{-b + \sqrt{b^2 - 4 a c}}{2 a}$と$\frac{-b - \sqrt{b^2 - 4 a c}}{2 a}$からなるタプルを返す関数として定義されている．
\begin{lstlisting}[language=egison,numbers=none]
def qF f x :=
  match coefficients f x as list mathExpr with
    | [$a_0, $a_1, $a_2] -> qF' a_2 a_1 a_0

def qF' a b c :=
  ( ((- b) + sqrt (b ^ 2 - 4 * a * c)) / 2 * a
  , ((- b) - sqrt (b ^ 2 - 4 * a * c)) / 2 * a )
\end{lstlisting}
しかし，上記の記述ではアルゴリズムっぽさがないため，今回は平方完成によりこの解の公式を求めるアルゴリズムを書いていく．
平方完成とは，$x^2 + b x + c$のような二次式を$(x + b / 2)^2 - b^2 / 4 + c$のように変形する操作のことである．
\begin{lstlisting}[language=egison]
def qF' a b c :=
  match (a, b, c) as (mathExpr, mathExpr, mathExpr) with
    | (#1, #0, _) -> (sqrt (- c), - sqrt (- c))
    | (#1, _, _) ->
      (2)#((- (b / 2)) + %1, (- (b / 2)) + %2)
        (withSymbols [x, y]
          qF (substitute [(x, y - b / 2)] (x ^ 2 + b * x + c)) y)
    | (_, _, _) -> qF' 1 (b / a) (c / a)
\end{lstlisting}
平方完成を使ってを二次方程式を解くアルゴリズムを記述すると\lstinline{qF'}は3つのマッチ節からなる関数となる．
3行目のマッチ節は，$x^2 + c = 0$のような形の二次方程式を処理する．
このマッチ節は，$\sqrt{c}$と$-\sqrt{c}$という結果を返す．
4-7行目のマッチ節は，$x^2 + b x + c = 0$のような二次の係数が$1$である形の二次方程式を処理する．
このマッチ節が平方完成の操作をするこの関数の本体である．
\lstinline{substitute}はEgisonのライブラリ関数である．
\lstinline{substitute}は第一引数で指定された代入を第二引数の式に対しておこなう．
ここでは第一引数で，\lstinline{(x, y - b / 2)}という代入が指定されているため，$x = y - b / 2$という代入を式$x^2 + b x + c$に対しておこなう．
同時に複数の代入を指定するために，\lstinline{substitute}の第一引数はリストを引数に取るようになっている．
$x = y - b / 2$という代入をするとこの方程式は3行目のマッチ節でマッチする$y$についての二次方程式に変形される．
そのため，\lstinline{qF}関数にこの変形された方程式を渡せば，$y$についてこの二次方程式の解を得られる．
この$y$についての解に$\frac{-b}{2}$を加えることにより，もとの$x$についての解を得ることができる．
\lstinline{withSymbols}式は，局所シンボルを生成する構文である．
第一引数のリストで指定されたシンボル（上記の場合，\lstinline{x}と\lstinline{y}）を第二引数の式の中でシンボルとして使うことができる．
\lstinline{withSymbols}式のこの機能のおかげで，もしプログラムの別の箇所で，\lstinline{x}や\lstinline{y}にたとえば具体的な整数が束縛されていたとしても，\lstinline{withSymbols}式の内部では，その影響を考えなくてよくなる．
8行目のマッチ節は，二次の係数が$1$でない二次方程式を処理する．
二次の係数で方程式を割ることにより，二次の係数が$1$の方程式に変換し，ふたたび\lstinline{qF}関数にわたす処理をしている．

二次方程式より一段と複雑になるが，三次方程式・四次方程式を解くプログラムを書いてみると面白い．
これらのプログラムはEgisonのソースコードの\texttt{sample/math/algebra}ディレクトリ以下に公開されている．

\section{微分するプログラム - 数式に対するパターンマッチ}\label{cas-derivative}

Egisonには微分計算をするための関数\lstinline{d/d}がライブラリ関数として実装されている．
この関数は第一引数の関数を第二引数のシンボルについて微分した結果を返す．
関数の名前を\lstinline{d/d}のようにすることで数式に近い形で微分を表現できている．
Egisonは\lstinline{/}を変数名に使うことを許している．
\begin{lstlisting}[language=egison,numbers=none]
d/d x x -- 1
d/d (x^2) x -- 2 * x
d/d (exp x) x -- exp x
d/d (log x) x -- 1 / x
d/d (x * log x) x -- log x + 1
d/d (1 / log x) x -- -1 / (x * (log x)^2)
\end{lstlisting}
本節はこの関数の実装を紹介する．

\lstinline{d/d}は以下のように定義される．
\begin{lstlisting}[language=egison]
def d/d f x :=
  match f as mathExpr with
    -- `シンボルの微分`
    | #x -> 1
    | ?isSymbol -> 0
    -- `関数適用の微分`
    | #'exp $g -> exp g * d/d g x
    | #'log $g -> 1 / g * d/d g x
    | #'sqrt $g -> 1 / (2 * sqrt g) * d/d g x
    | #'(^) $g $h -> f * d/d (log g * h) x
    | #'cos $g -> (- sin g) * d/d g x
    | #'sin $g -> cos g * d/d g x
    -- `定数の微分`
    | #0 -> 0
    | _ * #1 -> 0
    -- `項の微分`
    | #1 * $fx ^ $n -> n * fx ^ (n - 1) * d/d fx x
    | $a * $fx ^ $n * $r -> a * d/d (fx ^' n) x * r + a * fx ^' n * d/d r x
    -- `多項式の微分`
    | poly $ts -> sum (map 1#(d/d %1 x) ts)
    -- `商の微分`
    | $p1 / $p2 ->
      let p1' := d/d p1 x
          p2' := d/d p2 x
       in (p1' * p2 - p2' * p1) / p2 ^ 2
\end{lstlisting}
第一引数の$f$を数式としてパターンマッチしている．
4-5行目はシンボルの微分を定義している．
もし，\lstinline{f}が\lstinline{x}そのものであったら，$1$を返し，\lstinline{f}が\lstinline{x}以外のシンボルであったら$0$を返すように定義されている．
7-12行目はいくつかの基本的な関数の適用について合成関数の微分$\frac{df(g(x))}{dx} = \frac{df(g(x))}{dg(x)} \frac{dg(x)}{x}$を記述している．
たとえば，7行目は指数関数の微分をしている．
このマッチ節の本体は$\frac{de^{g(x)}}{dx} = e^{g(x)} \frac{dg(x)}{dx}$を表現している．
また，8行目は対数関数の微分をしている．
このマッチ節の本体は$\frac{dlog(g(x))}{dx} = \frac{1}{g(x)} \frac{dg(x)}{dx}$を表現している．
14-15行目は定数項に対する微分を定義している．
定数項は微分すると$0$になる．
17-18行目は，ここまでに微分を定義した項を掛け合わせた項についての微分を定義している．
$x^n$は$n x^{n-1}$に微分されることと，$f g$は$f' g + f g'$に微分されることを利用して定義している．
ここでも合成関数の微分の公式が使われている．
20行目は多項式の微分が定義されている．
ここまでで単項式についての微分が定義されているので，それらを足し合わせたものについては，それぞれの項について微分して足し合わせてやればよい．
22-24行目は多項式を分母と分子にもつ式に対する微分が定義されている．
商の微分の公式を使って定義されている．

\section{ベクトルや行列の計算 - 添字記法}\label{cas-linear}

Egisonの数式処理システムとしての特徴に，ベクトルや行列，そしてそれらを一般化したテンソルを扱う計算を簡潔に記述できるというものがある．
これは二十世紀初頭に微分幾何や相対性理論の研究の過程で発明されたテンソルの添字記法をプログラミング言語の機能としてEgisonがサポートしているためである．

添字記法は，テンソルに添字を付加することによって，さまざまな種類のテンソルの掛け算を表現する．
たとえば，ベクトル同士の掛け算には，テンソル積・アダマール積・内積の三種類があるが，これらは以下のように添字を使い分けることによって表現される．
\begin{lstlisting}[language=egison,numbers=none]
[| x1, x2 |]_i .[| y1, y2 |]_j -- [| [| x1 * y1, x1 * y2 |], [| x2 * y1, x2 * y2 |] |]_i_j
[| x1, x2 |]_i .[| y1, y2 |]_i -- [| x1 * y1, x2 * y2 |]_i
[| x1, x2 |]~i .[| y1, y2 |]_i -- x1 * y1 + x2 * y2
\end{lstlisting}
上記の例のように，Egisonではベクトルは成分を\lstinline{[| |]}で囲むことにより表現する．
また行列は，ベクトルのベクトルとして表現する．
関数型プログラミングの視点で添字記法は，テンソルの掛け算関数のパラメータとして，テンソルだけでなく添字も追加で渡すことによって一つの関数に複数の役割を持たしていると説明できる．

添字記法を言語機能としてもつプログラミング言語はほかにも存在する．
そのなかでEgisonの特徴は，関数型プログラミングとうまく調和するかたちで添字記法を組み込んでいるところである．
そのおかげで，高階関数と添字記法を組み合わせて使うことができる．
たとえば，$g_{i_{1}j_{1}} g_{i_{2}j_{2}} ... g_{i_{n}j_{n}}$のような数式を\lstinline{foldl}関数を使って以下のように記述することができる．
\begin{lstlisting}[language=egison,numbers=none]
foldl . 1 (map (\x -> g_[i_x]_[j_x]) [1..n])
\end{lstlisting}

添字記法はEgisonの重要な言語機能であるため，第\ref{tensor}章と第\ref{diffgeo}章で詳しく紹介する．

\section{数式の出力形式}

``\texttt{-M}''オプションで出力形式を指定することができる．
``\texttt{egison -M latex}''を実行すると，LaTeX形式で結果を出力する．
\begin{lstlisting}[language=egison,numbers=none]
$ egison -M latex
Egison Version 3.7.14 (C) 2011-2018 Satoshi Egi
https://www.egison.org
Welcome to Egison Interpreter!
> (+ (sqrt a) b)
#latex|\sqrt{a} + b|#
\end{lstlisting}
``\texttt{egison -M mathematica}''を実行すると，Mathematica形式で結果を出力する．
\begin{lstlisting}[language=egison,numbers=none]
$ egison -M mathematica
Egison Version 3.7.14 (C) 2011-2018 Satoshi Egi
https://www.egison.org
Welcome to Egison Interpreter!
> (+ (sqrt a) b)
#mathematica|Sqrt[a] + b|#
\end{lstlisting}

LaTeX出力はJupyter Notebookとも連携している．
Jupyter NotebookのEgisonプラグインを利用すれば，図\ref{fig:riemann-curavture-of-S2}，図\ref{fig:hodge-laplacian}のようにインタラクティブに計算結果が数式として表示される．

\begin{figure}[t]
  \centering
  \includegraphics[width=14.0cm,bb=0 0 1618 870]{jupyter-sample.png}
  \caption{Jupyter Notebook上でEgisonを実行する様子}
  \label{fig:cas-jupyter}
\end{figure}

